\begin{otherlanguage}{italian}
\chapter*{Sommario}
\addcontentsline{toc}{chapter}{Sommario}
I dispositivi IoT possono presentare comportamenti di rete caratteristici,
che consentono di costruire profili di normalità per il rilevamento delle
anomalie. Questa tesi presenta LM-IDNet, un rilevatore specifico per
dispositivo che utilizza i conteggi di pacchetti TCP, UDP ordinario, SSDP
e ARP in finestre complete di dieci minuti. Un modello
Dirichlet--multinomiale rappresenta la composizione del traffico ed è
affiancato da controlli sul volume totale e sui pacchetti ARP, calibrati
congiuntamente. La stima e la calibrazione utilizzano periodi distinti
di traffico presunto benigno; le etichette degli attacchi sono impiegate
nella valutazione fuori linea e nella selezione del metodo. Il sistema
integra il metodo numerico di Languasco--Migliardi con la stima di Minka,
aggiornamenti controllati di soglia e modello e uno stato persistente.

Gli esperimenti distinguono la valutazione su traffico etichettato di
Chromecast e di una videocamera Samsung dall'analisi numerica su un
carico D-Link privo di etichette di attacco. Una riproduzione conclusiva
con configurazione fissata comprende 1.008 finestre uniche etichettate:
la precisione aggregata è 82,35\%, il richiamo 77,78\% e F1 80,00\%.
F1 varia sensibilmente tra i dispositivi: 61,54\% e 90,91\%. Sono
rilevati otto dei nove intervalli di attacco interamente osservati;
per sette l'allarme è disponibile entro la fine dell'attacco. Le
etichette della partizione successiva erano già state esaminate in
analisi retrospettive, limitando le conclusioni sulla generalizzazione
indipendente. Le modalità conservative non accettano aggiornamenti
nella riproduzione conclusiva e producono le stesse decisioni della
modalità statica. Esperimenti storici e con durate diverse mostrano
che l'adattamento può peggiorare il rilevamento. Finestre più lunghe
possono migliorare le proprie metriche, ritardando però gli allarmi.

LM e SciPy mantengono tutte le decisioni di sviluppo verificate, con
risultati numerici vicini e tempi dipendenti dall'operazione. Una
simulazione Docker Linux/amd64, limitata a una CPU e 512~MiB, elabora
16.000 finestre con un tempo medio vicino a 30~ms e un massimo di
memoria di processo vicino a 69~MiB. La finestra più lenta richiede
circa 5,08~s durante una ristima poi rifiutata. Le misure riguardano
finestre già aggregate, escludendo acquisizione dal vivo e gateway
fisici. La tesi fornisce un prototipo riproducibile e interpretabile
e identifica i limiti della rappresentazione, dell'adattamento e
della valutazione da approfondire con studi indipendenti.
\end{otherlanguage}
